\section{Compatibility of the two beams}

The two beams address different questions and use different state semantics. This is an
advantage if the interface is chosen carefully.

\subsection{Objects that should not be identified}

Three tempting identifications should be rejected.

\paragraph{Knowledge-acquisition cost is not competence.}
Garicano's \(c\) is the marginal cost of acquiring knowledge. Gutjahr's \(z_{it}\) is a
competence score and \(\phi(z_{it})\) is productive efficiency. A relation such as
\(c_t=\psi(C_t)\) is not supplied by either theory and is unnecessary for joining them.

\paragraph{Knowledge set is not automatically proficiency.}
Garicano's
\[
A_t\subseteq\Q
\]
records which problem solutions are known. A continuous competence variable \(C_t(q)\)
can instead represent how effectively a known or learnable class of tasks is performed.
Defining
\[
A_t=\{q:C_t(q)\geq\vartheta\}
\]
would be an additional modeling assumption, not a consequence of the foundations reviewed
here.

\paragraph{Allocation, real work and learning exposure are distinct.}
An allocation decision determines what work is attempted and by whom. Real work is the
resource actually consumed in execution. Learning exposure is the part and type of the
resulting experience relevant to a competence transition. Keeping these objects distinct
allows failures, partial completion, feedback quality, support work, or non-learning work
without changing the conceptual skeleton.

\subsection{Compatible common structure}

Beam~1 ends with an operational allocation:
\[
(F_t,K_t,\kappa_t)\longrightarrow a_t.
\]
Execution maps allocation to real work:
\[
w_t=\mathcal{W}(a_t,F_t,\text{execution conditions}).
\]
Work and outcomes map to learning-relevant exposure:
\[
e_t=\mathcal{E}(w_t,\text{execution outcomes}),
\]
and Beam~2 uses that exposure in a competence transition:
\[
C_{t+1}=\mathcal{L}(C_t,e_t).
\]

The maps \(\mathcal W\), \(\mathcal E\), and \(\mathcal L\) are part of an HLS
instantiation; they are not supplied jointly by the source papers. The interface is therefore
a candidate modeling scaffold governed by the canonical HLS ontology, not an established
cross-paper identification.
